%% CAPITULO 1 -- INTRODUCAO

\section{Introdução}

As ouvidorias públicas recebem diariamente manifestações de cidadãos, entre reclamações, denúncias, solicitações e sugestões, que precisam ser classificadas e encaminhadas ao órgão competente antes de qualquer resposta. Esse trabalho de triagem consome tempo e depende do julgamento de servidores, o que torna o tratamento das demandas mais lento e menos uniforme à medida que o volume cresce.

Nesse cenário, técnicas de inteligência artificial e de processamento de linguagem natural têm sido propostas para apoiar a triagem de demandas dirigidas à administração pública. Estudos recentes, como os de \citeonline{buryakov2024} e \citeonline{majumder2025}, indicam que classificadores e modelos de linguagem podem atribuir demandas às unidades responsáveis, embora ainda restem dúvidas sobre a supervisão humana e a rastreabilidade das decisões.

\subsection{Objetivo geral}

Analisar, em caráter exploratório, como a inteligência artificial pode apoiar a triagem de manifestações em ouvidorias públicas, preservando a supervisão humana e a rastreabilidade das decisões tomadas pelo sistema.

\subsection{Objetivos específicos}

\begin{itemize}
  \item mapear abordagens de inteligência artificial usadas na triagem de demandas públicas;
  \item comparar o contexto institucional e a tarefa tratada em estudos de referência;
  \item delinear um fluxo de triagem com supervisão humana para ouvidorias públicas.
\end{itemize}

\subsection{Problema}

Que abordagens de inteligência artificial têm sido utilizadas para apoiar a triagem de manifestações dirigidas à administração pública, e que cuidados são necessários para empregá-las com supervisão humana e rastreabilidade das decisões?

\subsection{Justificativa}

O tema é relevante porque o aumento no volume de manifestações pressiona as ouvidorias, e decisões de encaminhamento equivocadas atrasam a resposta ao cidadão. Um esqueleto de pesquisa bem delimitado permite organizar a literatura e orientar etapas posteriores de investigação.
